\section{I/O}
\label{section:preparation:io}

\begin{assessmentcrime}
 The code writes a lot of information to the terminal. 
\end{assessmentcrime}

String operations are excessively expensive.



In this context, it is important that there's no I/O.
We don't want to study complex outputs; we are busy enough with handling the output of performance analysis tools.

\begin{assessmentcrime}
 The code reads or write to files excessively. 
\end{assessmentcrime}

\noindent
More importantly, if our code writes excessive amounts of data, we
    almost certainly profile the I/O efficiency rather than the runtime
    behaviour. Unless we want to explicitly study I/O, we are well-advised to
    switch off any file outputs, database outputs, but also excessive file
    reads.
 A log file alone can become problematic when these files become suddenly huge or are updated frequently.
 In the best case, a performance assessment benchmark writes and reads barely any data at all, i.e.~it largely hard-coded and silent.


